\documentclass[10pt]{article}
\usepackage[margin=0.7in]{geometry}
\usepackage{amsmath}
\usepackage{booktabs}
\usepackage{siunitx}
\usepackage{xcolor}
\usepackage{enumitem}
\usepackage{helvet}
\renewcommand{\familydefault}{\sfdefault}
\setlength{\parindent}{0pt}
\setlist[itemize]{leftmargin=1.4em,itemsep=0.35em,topsep=0.3em}
\sisetup{detect-all}

\definecolor{accent}{RGB}{0,114,178}

\begin{document}

\begin{center}
{\Large \textbf{STAR PXL Single-Track Pointing Resolution Estimate}}\\[0.3em]
{\color{accent}\rule{0.92\linewidth}{0.8pt}}
\end{center}

\textbf{Detector inputs}
\begin{itemize}
  \item Two pixel layers at radii
  \[
    r_1=\SI{2.8}{cm}, \qquad r_2=\SI{8.0}{cm}.
  \]
  \item Pixel pitch: \(\SI{20.7}{\micro m}\times\SI{20.7}{\micro m}\), so the single-hit resolution is
  \[
    \sigma_{\rm hit}=\frac{20.7}{\sqrt{12}}=\SI{5.98}{\micro m}.
  \]
  \item Pixel material: \(0.52\%\,X_0\) per layer.
  \item Beam pipe: \(r_{\rm BP}=\SI{2.0}{cm}\), material \(0.24\%\,X_0\).
\end{itemize}

\textbf{Hit-resolution limit}
\begin{itemize}
  \item For a straight track \(y(r)=a+br\), the intercept at the collision point is
  \[
    a = \frac{r_2 y_1-r_1 y_2}{r_2-r_1}.
  \]
  \item Propagating the two independent hit errors gives
  \[
    \sigma_{a,{\rm hit}}
    =
    \sigma_{\rm hit}
    \sqrt{
      \left(\frac{r_2}{r_2-r_1}\right)^2+
      \left(\frac{r_1}{r_2-r_1}\right)^2
    }
    =
    \SI{9.74}{\micro m}.
  \]
\end{itemize}

\textbf{Multiple-scattering term}
\begin{itemize}
  \item Using the Highland approximation for \(\beta\simeq 1\),
  \[
    \theta_0
    \simeq
    \frac{13.6~{\rm MeV}}{p}
    \sqrt{\frac{x}{X_0}}
    \left[1+0.038\ln\left(\frac{x}{X_0}\right)\right].
  \]
  \item The beam pipe contributes
  \[
    \SI{2.0}{cm}\times\frac{0.5135~{\rm mrad}}{p}
    =
    \frac{\SI{10.27}{\micro m\,GeV}/c}{p}.
  \]
  \item The first PXL layer contributes
  \[
    \SI{2.8}{cm}\times\frac{0.7847~{\rm mrad}}{p}
    =
    \frac{\SI{21.97}{\micro m\,GeV}/c}{p}.
  \]
  \item The outer PXL layer is the last measured point in this two-hit estimate, so its post-hit scattering does not affect the fitted pointing.
\end{itemize}

\textbf{Final fast estimate}
\[
  \sigma_{\rm point}(p)
  \simeq
  \sqrt{
    (\SI{9.74}{\micro m})^2
    +
    \left(\frac{\SI{24.25}{\micro m\,GeV}/c}{p}\right)^2
  },
  \qquad p~{\rm in~GeV}/c.
\]

\begin{center}
\begin{tabular}{@{}S[table-format=2.1] S[table-format=3.1]@{}}
\toprule
{\(p\) [GeV/\(c\)]} & {\(\sigma_{\rm point}\) [\(\si{\micro m}\)]} \\
\midrule
0.2 & 121.7 \\
0.5 & 49.5 \\
1.0 & 26.1 \\
2.0 & 15.6 \\
5.0 & 10.9 \\
10.0 & 10.0 \\
\bottomrule
\end{tabular}
\end{center}

\vspace{0.5em}
This estimate is for central tracks and uses the same resolution in \(r\phi\) and \(z\).

\end{document}
